\bookchapter{第二十二章：AIはスマホでどう働くのか}

\chapterintro{二、音声アシスタント：あなたの「声の秘書」はどうやってあなたを理解する？}

\booksection{聞き取る（音声認識）}

：音波をテキストに変換。

\booksection{理解する（自然言語理解）}

\begin{enumerate}
  \item \textbf{：意図を理解、キーエンティティを抽出。}実行と回答
  \item \textbf{：対応するAppやサービスを呼び出し、音声合成で「話して」聞かせる。}三、オンデバイスAI vs. クラウドAI
  \item \textbf{オンデバイスAIは計算をスマホ内に留め、速度が速く、プライバシーが高く、オフラインで利用可能。クラウドAIはネット接続が必要だが、より複雑なタスクを処理できる。}四、AI大規模モデルとスマホ：あなたの常時携帯スマートアシスタント
\end{enumerate}

\booksection{2024-2025年、AI大規模モデル（DeepSeek、ChatGPT、文心一言、通義千問など）の爆発がスマホの使い方を深く変えました。スマホはもはや指示を実行するだけのツールではなく、理解、推論、創作できるスマートパートナー。}

AI大規模モデルとは？

\booksection{AI大規模モデル（Large Language Model、LLM）は膨大なテキストデータで訓練されたAIシステムで、自然言語を理解·生成できます。インターネット上のほぼすべての文字を読んだ「スーパー優等生」と想像でき、質問に答え、文章を書き、翻訳し、プログラミングし、データ分析できる。}

2025-2026年、AI 大規模モデル（DeepSeek R1、ChatGPT-5、文心一言5.0、通義千問3 など）の爆発的普及はスマホの使い方を根本的に変えました。スマホはもはや命令を実行するだけのツールではなく、理解し、推論し、創作するインテリジェントなパートナーです。オンデバイス AI チップ（NPU 演算力80 TOPS 突破）により、大規模モデル推論が初めてスマホ上でローカル実行可能になり、インターネット接続なしでも AI が使えます。

\booksubsection{対話アシスタント}

：DeepSeek、ChatGPT、豆包、KimiなどのAppで、いつでも質問、翻訳、文案作成、記事要約が可能。伝統的な検索エンジンよりスマート——リンクだけでなく、直接答えをくれる。

\booksubsection{執筆と創作}

\begin{itemize}
  \item \textbf{：AIにメール作成、会議議事録要約、文章推敲、クリエイティブ文案生成を依頼。}画像生成と編集
  \item \textbf{：AIでワンタップ写真の通行人を消去、画像境界を拡張、スタイル変換、さらにはテキスト記述から画像を生成。}リアルタイム翻訳
  \item \textbf{：通話リアルタイム翻訳、対面会話翻訳、写真翻訳、海外旅行で言語の壁なし。}学習補習
  \item \textbf{リアルタイム翻訳}オンデバイス大規模モデル：スマホ自身でAIを実行
  \item \textbf{最新のフラッグシップスマホはローカルで小型AI大規模モデル（通常7Bパラメータ程度）を実行可能、もたらす変化は：}ネットワーク不要
\end{itemize}

\booksubsection{：ネット切断でもAI機能使用可能、機内、地下鉄でもOK。}

プライバシーがより安全

\begin{itemize}
  \item \textbf{：あなたの対話とデータはクラウドにアップロードされず、完全にスマホローカルで処理。}応答がより速い
  \item \textbf{：ネットワーク遅延なし、AI応答はほぼ瞬時。}現在の制限
  \item \textbf{：オンデバイスモデルの能力はまだクラウド大規模モデルに及ばず、複雑なタスクは依然としてネット接続が必要。ただしチップのAI演算力向上に伴い、オンデバイスモデルはますます強力に。}システムレベルAI統合：
  \item \textbf{Apple Apple Intelligence}：iOS 18で導入されたAIシステム、執筆ツール、画像マジックワンド、Siri強化、通知要約などの機能を統合。オンデバイス処理を優先、複雑なタスクのみクラウドへ。
\end{itemize}

\booksubsection{Huawei盤古大規模モデル}

\begin{itemize}
  \item \textbf{：HarmonyOSに統合、小芸アシスタントが大規模モデル能力を獲得、ドキュメント要約、スマート画面認識、AIレタッチなどをサポート。}Xiaomi澎湃AI
  \item \textbf{：システムレベルAI能力、AI画像検索、AI要約、AIリアルタイム字幕、オンデバイス大規模モデル対話などをサポート。}Samsung Galaxy AI
  \item \textbf{：通話リアルタイム翻訳、ノートAI要約、画像生成編集、サークル検索などの機能。}五、いくつかの典型的な応用シーンと実践
  \item \textbf{スマホAIの「超能力」をより直感的に感じてもらうため、以下5つの具体的なシーンで、AIが日常生活でどう「助け舟を出す」かをお見せします：}シーン 1：リアルタイム翻訳——海外旅行の常時携帯通訳
\end{itemize}

\booksection{実際の表現}

：パリのレストランでフランス語が分からない？iPhoneの「翻訳」Appは対面双方向会話モードをサポート、双方が各自の言語の字幕を見られ、同時通訳も可能。Android陣営のSamsung Galaxy AI通話リアルタイム翻訳は、電話で直接中国語を話し、相手には英語が聞こえ、その逆も可能。

\booksubsection{背後の原理}

\textbf{：コンピュータビジョン（CV）モデルが画像内容をベクトル化（Embedding）、テキストクエリと類似度マッチング。この「マルチモーダル検索」はスマホアルバム管理の大きな飛躍。}シーン 2：AI消去——通行人に「破壊された」観光写真を救済

\textbf{実際の表現}：観光地で絶景の自撮りを撮ったが、背景に他の観光客がぎっしり？「マジック消去」で通行人を囲むと、スマホが自動認識して通行人を「塗り消し」、さらに背景をスマートに埋め合わせ、効果はしばしば驚異的。

\booksubsection{背後の原理}

\textbf{：コンピュータビジョン（CV）モデルが画像内容をベクトル化（Embedding）、テキストクエリと類似度マッチング。この「マルチモーダル検索」はスマホアルバム管理の大きな飛躍。}シーン 3：AIレタッチ——ワンタップで「撮影マスター」

\textbf{実際の表現}：曇りの日の灰色っぽい風景写真を撮った？AI「マジックスカイ」でワンタップで空を晴れ、夕焼け、さらには星空に交換。AI画像拡張で縦構図を横構図にでき、スマホが画面外のシーンを「推測」。

\booksubsection{背後の原理}

\textbf{：コンピュータビジョン（CV）モデルが画像内容をベクトル化（Embedding）、テキストクエリと類似度マッチング。この「マルチモーダル検索」はスマホアルバム管理の大きな飛躍。}シーン 4：ドキュメント処理——オフィスワーカーの効率倍増器

\textbf{実際の表現}：会議終了後、会議録音をAIに渡すと、自動的に構造化された会議議事録を整理、ToDo項目を抽出、要約を生成。10ページの英語PDFレポートをAIに渡せば、10秒で核心結論を抽出。

\booksubsection{背後の原理}

\textbf{：コンピュータビジョン（CV）モデルが画像内容をベクトル化（Embedding）、テキストクエリと類似度マッチング。この「マルチモーダル検索」はスマホアルバム管理の大きな飛躍。}シーン 5：AI画像検索と物体認識——あなたの「視覚検索エンジン」

\textbf{実際の表現}：街で知らない花を見かけ、カメラを起動するとAIが即座に「これはムクゲです」と教える。アルバムで去年海辺で撮った写真を探す時、「海」と検索すると、AIが海を含むすべての写真を自動検出、手動で探す必要なし。

\booksubsection{背後の原理}

\textbf{：コンピュータビジョン（CV）モデルが画像内容をベクトル化（Embedding）、テキストクエリと類似度マッチング。この「マルチモーダル検索」はスマホアルバム管理の大きな飛躍。}六、プロンプト最適化：AIをあなたの指揮下に置く「コミュニケーション公式」

\textbf{実際の表現}最も実用的で最も使いやすい「役割-背景-タスク-制約」コミュニケーション思路

\booksection{を紹介し、大量の実例を交える。}

なぜフレームワークが必要？\textbf{多くの人がAIをBaiduのように使う——キーワードをいくつか入力する。しかしAIは検索エンジンではなく、「読解してから回答する」言語モデル。}より多くの文脈を与え、何を欲しているかを明確にし、脱線を制限する

\begin{tcbinfo}{と、精確な結果をくれる。フレームワークはあなたの頭の中の「良い答えとは何か」をAIが理解できる指示に翻訳するもの。}
1. 役割（Role）：AIに誰であるかを伝える\textbf{AIに具体的な「アイデンティティ」を与えると、対応分野の専門知識と表現方法を呼び出す。}役割設定が具体的なほど、出力が専門的

\end{tcbinfo}

\booksubsection{。}

AIに具体的な「役割」を与えると、その分野の専門知識と表現スタイルを引き出してくれます。\textbf{「あなたは}。

\booksubsection{で。」（長さ＋形式）}

\begin{itemize}
  \item なぜ有効？\textbf{複雑な概念をシンプルに説明するのが得意な中学物理の先生}です。」（科学普及問題）
  \item なぜ有効？\textbf{審美眼があり、多くを読み、独自のスタイルを持つ図書編集者}です。」（執筆問題）
  \item なぜ有効？\textbf{毎日会社員に朝食を作る料理ブロガー}です。」（生活問題）
  \item なぜ有効？\textbf{忍耐強く、穏やかで、初心者を決して嘲笑しない中医師}です。」（健康問題）
  \item なぜ有効？\textbf{役割はAIの}知識範囲
\end{itemize}

\textbf{良いプロンプトを書くのは一回きりのことではなく、}表現スタイル\textbf{、}、\textbf{を設定。役割のないAIは「何でもできるが何も深くない」ゼネラリスト。役割があれば、そのアイデンティティに入り、その専門家のように思考できる。}3. タスク（Task）：AIに何をするか明確に伝える\textbf{「あなたは誰、現状は何、制約条件は何」を事前に明確にすると、AIは}本当にあなたに適した

\booksubsection{答えを出せ、漠然とした回答ではなくなる。}

「自分は誰か、今どんな状況か、どんな制約があるか」を先にはっきり伝えると、AIは\textbf{「私のスマホは}Xiaomi 14、買って半年、最近ゲーム中に発熱がひどい

\booksubsection{で。」（長さ＋形式）}

\begin{itemize}
  \item 「私は\textbf{60歳の退職したおばさんで、目があまり良くなく、スマホをあまり使いこなせません}。」（ユーザー像）
  \item 「私は\textbf{2026年6月の英語六級試験を準備中、基礎は普通}です。」（学習背景）
  \item 「私は\textbf{三都市の小さなカフェを経営、月商3万、二人の従業員を雇用}しています。」（ビジネス背景）
  \item 「私は\textbf{出産後家で産後休養中、睡眠が著しく不足}しています。」（生活段階）
  \item なぜ有効？\textbf{背景情報はAIに}「誰のためにサービスするか」
\end{itemize}

\textbf{良いプロンプトを書くのは一回きりのことではなく、}「どんな条件下か」\textbf{を知らせる。同じ問題でも、学生と親には説明が異なり、都市人と地方人にはアドバイスが異なる。背景が欠けると、AIは「最大公約数」の答えしか出せず、しばしばピンとこない。}3. タスク（Task）：AIに何をするか明確に伝える\textbf{このステップは}動詞で始める

\booksubsection{（書く、列挙、比較、翻訳、説明、要約……）、タスクを}

具体的で実行可能\textbf{に記述。曖昧なタスクは曖昧な答えしか生まない。}（書く、列挙、比較、翻訳、説明、要約……）。タスクは\textbf{「}私がすぐに閉じるべき5つの権限を列挙

\booksubsection{で。」（長さ＋形式）}

\begin{itemize}
  \item 「\textbf{以下の英語メールを中国語に翻訳}し、ビジネス口調を維持してください。」（タスク具体）
  \item 「\textbf{この4機種の3000元帯のスマホを比較}し、長所短所比較表を作成してください。」（形式明確）
  \item 「\textbf{小学生にも分かる言葉で「5Gとは何か」を説明}してください。」（スタイル明確）
  \item 「\textbf{私の朝食店に3つの短動画文案を執筆}し、それぞれ会社員、カップル、子供連れの親向けにしてください。」（セグメント明確）
  \item 「\textbf{タスク部分はAIの行動の「指令中枢」。}動詞が明確で要求が具体的なほど、AIは「何を納品するか」を把握
\end{itemize}

\textbf{良いプロンプトを書くのは一回きりのことではなく、}4. 制約（Constraints）：AIに何をしないかを伝える\textbf{「制約」は「拘束」のように聞こえるが、実は}AIをあなたに最も適した出力に収束させる

\booksubsection{強力なツール。一般的な次元は：文字数、形式、トーン、回避項目、必須項目。}

「制限」は「縛り」のように聞こえますが、実際には\textbf{「}200字以内、箇条書き

\booksubsection{で。」（長さ＋形式）}

\begin{itemize}
  \item 「\textbf{トーンは親しみやすく自然に、友人との会話のように}、『まず、次に、さらに』のような役所語は使わないで。」（トーン）
  \item 「\textbf{専門用語は使わず}、専門用語に出会ったらついでに説明して。」（用語制限）
  \item 「\textbf{各段落は3行以内}、スマホで読みやすく。」（レイアウト）
  \item 「\textbf{必ず含める：使用シーン、長所短所、私の予算との適合度}。」（必須項目）
  \item 「\textbf{具体的な投資アドバイスはしない}、原理だけ説明して。」（境界）
  \item 「\textbf{制約がないと、AIは「自由発揮」しやすい——学術的すぎたり、長すぎたり、脱線したり。明確な制約＝明確な「納品基準」、双方の期待が揃い、結果がコントロール可能に。}5. プロンプト品質評価：良いプロンプトの5つの特徴
\end{itemize}

\textbf{良いプロンプトを書くのは一回きりのことではなく、}書いたらこの5項目を自己チェック

\booksubsection{してからAIに送信：}

具体的で抽象的でない\textbf{：「記事を書いて」ではなく「800字、3セクション、初心者ママ向けのスマホ選びガイドを書いて」。}背景を含み孤立していない

\begin{enumerate}
  \item \textbf{：AIに「私は誰、なぜこれを聞くか」を伝えれば、「あなた用にカスタマイズ」された答えが出せる。}明確な出力形式
  \item \textbf{：表、リスト、段落、それともJSON？事前に言う。}制約が適切
  \item \textbf{：文字数、スタイル、回避項目、3つのうち1つでも明確にすれば顕著な向上。}AIが検証可能
  \item \textbf{：聞き終わった後、AIにあなたの要求を復唱させる；はっきり言えなければ、はっきり伝わっていない。}6. よくある問題と解決法
  \item \textbf{症状}解決法
\end{enumerate}

\booksubsection{回答が漠然、実用性がない}

\begin{center}
\begin{tabularx}{\linewidth}{|X|X|}
\hline
\textbf{背景を追加（あなたは誰/どんなシーン）} & \textbf{回答が学術的/公式すぎる} \tabularnewline
\hline
トーン制約を追加＋「役所語」拒否リスト & 回答が脱線/質問に答えていない \tabularnewline
タスクを動詞で言い直し、「これだけ答えて」を要求 & 回答が長すぎて読めない \tabularnewline
文字数を硬く制限（200字以内） & AIがもっともらしく嘘をつく \tabularnewline
「不確かな場合は『分かりません』と表示」を追加＋重要事実の自己チェック & 第一版が不満足 \tabularnewline
書き直さず、元のプロンプト後に「もう少し…して」を追加 & 7. 3つの完全なシナリオ事例 \tabularnewline
シナリオ1：生活アシスタント（朝食組み合わせ提案） & プロンプト \tabularnewline
\hline
\end{tabularx}
\end{center}

\booksubsection{：}

\begin{tcbinfo}{あなたは栄養バランスに優れた朝食ブロガーです（}
\textbf{シニア行政アシスタント}：

背景\textbf{）。}7日間飽きない朝食の組み合わせをリストアップしてください\textbf{（}タスク\textbf{）。各朝食の要件：(1) 準備時間15分以内 (2) 栄養バランス（タンパク質＋炭水化物＋野菜·果物）(3) 老若男女向け (4) 一般的なスーパーで買える食材（}制限\textbf{）。}）。各朝食の要件：(1)15分以内で作れる (2)栄養バランス（タンパク質＋炭水化物＋野菜）(3)老若問わず食べられる (4)食材は普通のスーパーで買える（\textbf{：役割が専門性を定義し、背景が3世代の共通点を示し、タスクが明確な成果物を指定し、4つの制限が実行可能性を保証します。}ログイン後、右上の「+」をクリックし、「New repository」を選択。

\textbf{プロンプト}：

\end{tcbinfo}

\begin{tcbinfo}{あなたは}
\textbf{シニア行政アシスタント}：

Q2マーケティング推進案を議論しました\textbf{（背景）。}実行可能な12週間の試験対策計画を立ててください\textbf{（タスク）。要件：(1) 毎週のリスニング·スピーキング·リーディング·ライティングの時間配分を明確に (2) 毎月の重点突破方向を提示 (3) 適切な復習資料を3冊推奨（選定理由も記載） (4) 表形式で出力（制限）。}（背景）。\textbf{：役割が専門的背景を定義し、背景がスコアと時間のアンカーを示し、タスクが「段階別計画」を明確にし、制限が「表化」と「実行可能性」を要求します。}シナリオ3：仕事効率化（議事録整理）

\textbf{プロンプト}：

\end{tcbinfo}

\begin{tcbinfo}{あなたは}
\textbf{シニア行政アシスタント}：

Q2マーケティング推進案を議論しました\textbf{（背景）。}以下の会議記録を構造化された議事録に整理してください\textbf{（タスク）。要件：(1) 「決議事項／ToDo事項／リスクと異論」の3部構成 (2) ToDo事項には担当者と期限を明記（言及がなければ「要確認」と記載） (3) 総文字数は600字以内 (4) Markdown形式で出力（制限）。}（背景）。\textbf{：役割が行政の視点を定義し、背景が会議テーマを示し、タスクが構造化出力を明確にし、4つの制限が結果をそのまま再利用可能にします。}応用編：AIをあなたの「専属アシスタント」にする

\textbf{プロンプト}複数回対話の最適化

\end{tcbinfo}

\textbf{：AIの回答に不満がある時、「これは専門的すぎるので、小学生にもわかる言葉で説明し直してください」と言えます。}\par

思考の連鎖（CoT）

\begin{itemize}
  \item \textbf{：AIに「段階的に思考させる」と、回答の質が通常高くなります：「この数学の問題を段階的に分析してください」。}複数AI連携
  \item \textbf{：Aで初稿を書き、Bで推敲し、Cで翻訳する——組み合わせて使うと効果的です。}システムプロンプト
  \item \textbf{：AIの設定で自分の身分、好み、タブーなどを事前設定し、回答を自分のニーズに合わせます。}七、AIの限界と課題：清醒な判断力を保つ
  \item \textbf{チャレンジ1：複雑な指示}「ハルシネーション」問題
\end{itemize}

\booksection{：AIがもっともらしくでたらめを言うことがあります（専門用語でHallucination）。存在しない歴史的事件、誤った引用、虚偽の統計データをでっち上げます。重要な情報は必ず人工で確認してください。}

知識の期限

\begin{itemize}
  \item \textbf{：ほとんどのAIの訓練データには期限があり（例：2024年初頭）、それ以降の出来事を知りません。}推論能力の限界
  \item \textbf{：複雑な多段推論、数学的証明、論理的罠において、AIは依然として間違いを犯しやすいです。}プライバシーリスク
  \item \textbf{：クラウドAIサービスでは、あなたの会話がモデル訓練に使用される可能性があります。機密情報には端末側AIを優先使用してください。}偏見と価値観
  \item \textbf{：AIは訓練データの偏見を持つ可能性があり、独立した思考を保つ必要があります。}AI使用のヒント
  \item \textbf{AI生成の内容は不正確な場合があります（「ハルシネーション」問題）、重要な情報は必ず確認してください。}AIに個人情報（身分証番号、パスワード、銀行情報など）を入力しないでください。
\end{itemize}

\begin{tcbinfo}{プロンプトを活用：より具体的で明確な記述ほど、AIの回答の質が高くなります。}
\begin{itemize}
  \item AIはツールであり権威ではありません、独立した思考と判断力を保ってください。
  \item 重要な意思決定（医療、法律、財務）では、AIに頼らず専門家に相談してください。
  \item 応用編：AIをあなたの「専属アシスタント」にする
  \item これらの応用使い方を試して、AIの真の実力を体験してください：
  \item 複数回対話の最適化
\end{itemize}

\end{tcbinfo}

\begin{tcbpractice}{チャレンジ2：AI画像編集}
\begin{enumerate}
  \item \textbf{：通行人が写った観光写真を撮り、「魔法消去」や「AI消しゴム」機能を使って通行人をきれいに消去してください。AIが背景をどう「補完」するか観察してください。}チャレンジ3：プロンプト最適化
  \item \textbf{：まず「旅行ガイドを書いて」という簡単なプロンプトでAIに聞き、次にCRISPEフレームワークで詳細なプロンプトを書いてください。2回の回答の質の違いを比較してください。}第二十三章：プログラミング的思考とオートメーション入門
  \item \textbf{前二章では、スマートフォンに安全な「境界壁」を築き、壁の中で日に日に賢くなる「AIの頭脳」を見てきました。ここで、スマートフォンの究極のポテンシャルの一つを解放しましょう：プログラマーのように「考え」、設定した連続タスクを自動実行させること。「プログラミングなんてできないし、自分には関係ない」と思うかもしれません。安心してください、本章では一行のコードも書く必要はありません——現代のスマートフォンは「プログラミング」のハードルを、ドラッグ＆ドロップで済むレベルまで下げています。}一、オートメーションとは何か？
\end{enumerate}

\end{tcbpractice}

\begin{tcbinfo}{事例2：機種変更データ移行の完全フロー}
\begin{itemize}
  \item \textbf{シーン}: iPhone 13 Xiaomi 15 のクロスブランド機種変更。
  \item \textbf{手順}: (1) 旧機で iCloud/WeChat バックアップ (2) Apple ID ログアウト +「検索」をオフ (3) 新機で「Xiaomi 機種変更」App でQRスキャン (4) 移行項目を選択 (5) WeChat は独自の「チャット履歴移行」を使用 (6) 重要データ確認（連絡先/アルバム/WeChat） (7) 旧機を暗号化消去。
  \item \textbf{期待される結果}: データ損失ゼロ、旧機は安全に売却可能。
\end{itemize}

\end{tcbinfo}
